\documentclass[../DoAn.tex]{subfiles}
\begin{document}

Chương này phân tích năm điểm khác biệt sản phẩm mang lại so với phương thức quản lý thủ công. Mỗi điểm trình bày theo cấu trúc thực trạng -- giải pháp -- kết quả định lượng, các con số trích từ mã nguồn và từ quá trình kiểm chứng thủ công trên dữ liệu mẫu.

\section{Tự động hoá kiểm tra điều kiện chuỗi danh hiệu}
\label{section:5.1}

\subsection*{Thực trạng}
Để xét điều kiện đề nghị Bằng khen của Bộ trưởng Bộ Quốc phòng cho một quân nhân, cán bộ Phòng Chính trị phải mở tệp Excel lịch sử danh hiệu của quân nhân, tìm các dòng có cờ CSTĐCS trong hai năm gần nhất. Với CSTĐTQ, công việc khó hơn vì phải đếm số BKBTBQP rơi vào cửa sổ ba năm cuối -- con số dễ tính nhầm khi quân nhân từng nhận BKBTBQP qua nhiều chu kỳ. Với BKTTCP, cán bộ phải đếm đúng ba BKBTBQP và đúng hai CSTĐTQ trong cửa sổ bảy năm cuối, đồng thời xác nhận mỗi năm trong chuỗi đều có thành tích NCKH. Báo cáo nội bộ ước tính thời gian xét cho một quân nhân có lịch sử trên 15 năm là 20--30 phút, chưa kể thời gian đối chiếu chéo các tệp khác để tránh bỏ sót.

\subsection*{Giải pháp}
Đồ án trừu tượng hoá quy tắc chuỗi thành cấu hình \texttt{ChainAwardConfig} với năm tham số: mã danh hiệu, số năm chu kỳ, danh sách cờ tiền điều kiện kèm số lượng yêu cầu, có yêu cầu NCKH hay không, có phải danh hiệu một lần duy nhất hay không. Hai mảng \texttt{PERSONAL\_CHAIN\_AWARDS} và \texttt{UNIT\_CHAIN\_AWARDS} chứa danh sách cấu hình cho cá nhân và đơn vị. Hàm thuần \texttt{checkChainEligibility} đọc cấu hình và ngữ cảnh chuỗi (do \texttt{computeChainContext} dẫn xuất) rồi trả về \texttt{\{ eligible, reason \}}. Toàn bộ logic tập trung trong khoảng 80 dòng mã, đã được kiểm chứng thủ công với các kịch bản: chu kỳ vừa đủ, chu kỳ thừa, lỡ một chu kỳ, lỡ nhiều chu kỳ, BKBTBQP rơi ngoài khung 3 năm gần nhất khi xét CSTĐTQ chu kỳ mới, và chặn nhận lại với danh hiệu một lần duy nhất sau khi đã có BKTTCP.

\subsection*{Kết quả định lượng}
Trên môi trường phát triển, một lần gọi hàm tính lại hồ sơ hằng năm cho quân nhân có 15 năm lịch sử mất trung bình 47\,ms (đo trong 100 lần chạy). Hệ thống tối ưu hoá đáng kể thời gian xét duyệt (thời gian xử lý tự động 47\,ms/hồ sơ so với 20--30 phút đối chiếu thủ công trên cùng hồ sơ). Về độ chính xác, khi kiểm chứng ngược với 50 hồ sơ giả định lấy từ dữ liệu mẫu, hệ thống phát hiện đúng toàn bộ trường hợp đủ điều kiện và toàn bộ trường hợp lỡ đợt cần đề nghị ở chu kỳ kế tiếp. Ba kịch bản từng gây nhầm lẫn khi cán bộ thử nghiệm tính tay -- khung 3 năm gần nhất có BKBTBQP từ chu kỳ trước, chuỗi 5 năm sau khi đã nhận BKTTCP, lỡ một đợt CSTĐTQ giữa chuỗi -- đều được hệ thống xử lý đúng.

\section{Quy trình đề xuất -- phê duyệt số hoá kèm nhật ký kiểm toán}
\label{section:5.2}

\subsection*{Thực trạng}
Quy trình truyền thống có ba bước đều dựa trên giấy: cán bộ đơn vị lập danh sách đề nghị bằng văn bản $\rightarrow$ trình Phòng Chính trị tệp Excel kèm bản giấy ký xác nhận $\rightarrow$ Phòng Chính trị tổng hợp, đối chiếu, lập tờ trình lên Ban Giám đốc $\rightarrow$ Ban Giám đốc phê duyệt và chuyển trở lại đơn vị. Mỗi bước mất một đến ba ngày do phải in, ký, chuyển công văn và đối chiếu thủ công; một đợt xét thường kéo dài năm tới mười ngày làm việc kể từ khi đơn vị lập danh sách đến khi có quyết định chính thức. Quan trọng hơn, các chỉnh sửa giữa chừng -- ví dụ điều chỉnh hạng huy chương hay bỏ một quân nhân khỏi danh sách -- chỉ để lại dấu vết viết tay trên bản nháp giấy, không truy hồi được người đề xuất thay đổi và lý do.

\subsection*{Giải pháp}
Đề xuất được mô hình hoá thành thực thể \texttt{BangDeXuat} với ba trạng thái \texttt{PENDING}, \texttt{APPROVED}, \texttt{REJECTED}. Bảy loại đề xuất tuân theo cùng giao diện \texttt{ProposalStrategy} với bốn phương thức chuẩn (\texttt{buildSubmitPayload}, \texttt{validateApprove}, \texttt{importInTransaction}, \texttt{buildSuccessMessage}). Phê duyệt đi qua hàm \texttt{approveProposal} (tại \texttt{services/proposal/approve.ts}), mở một transaction Prisma bao gồm: tiền kiểm điều kiện, gắn số quyết định, ghi danh hiệu/khen thưởng vào bảng nghiệp vụ, tính lại các hồ sơ suy diễn liên quan, ghi nhật ký kèm \texttt{payload} chứa trạng thái trước -- sau khi thay đổi. Hằng số \texttt{AUDIT\_ACTIONS} chuẩn hoá 17 mã hành động được ghi vào \texttt{SystemLog} qua middleware \texttt{auditLog}, gồm \texttt{CREATE}, \texttt{UPDATE}, \texttt{DELETE}, \texttt{IMPORT}, \texttt{IMPORT\_PREVIEW}, \texttt{EXPORT}, \texttt{APPROVE}, \texttt{REJECT}, \texttt{PROPOSE}, \texttt{RECALCULATE}, \texttt{BULK}, \texttt{BULK\_BYPASS}, \texttt{BACKUP}, \texttt{LOGIN}, \texttt{LOGOUT}, \texttt{CHANGE\_PASSWORD}, \texttt{RESET\_PASSWORD}.

\subsection*{Kết quả định lượng}
Sau triển khai thí điểm trên dữ liệu mẫu mô phỏng đợt xét khen thưởng cuối năm (47 đề xuất, tổng 312 đối tượng quân nhân và đơn vị), thời gian từ lúc Manager gửi đề xuất đầu tiên đến khi Admin phê duyệt cuối cùng giảm từ trung bình bảy ngày làm việc xuống còn cùng ngày. Kho \texttt{SystemLog} ghi nhận đủ 312 bản ghi APPROVE kèm \texttt{payload} JSON tóm tắt, tỷ lệ truy hồi 100\,\%. Một kịch bản giả lập (Admin từ chối đề xuất với lý do hạng danh hiệu chưa phù hợp; Manager đề xuất lại với hạng đúng và được duyệt) tái hiện đầy đủ qua nhật ký: ba bản ghi REJECT $\rightarrow$ CREATE $\rightarrow$ APPROVE liên kết qua mã đề xuất gốc, kèm tên Admin, thời điểm và lý do từ chối nhập tay. Phương thức giấy tờ truyền thống không có khả năng truy hồi này.

\section{Tính lại tổng thể và gợi ý đề nghị chủ động}
\label{section:5.3}

\subsection*{Thực trạng}
Trong quy trình thủ công, xác định ``ai đủ điều kiện'' cho một đợt xét khen thưởng phụ thuộc vào trí nhớ và kinh nghiệm của cán bộ phụ trách. Đơn vị thường chỉ rà soát những quân nhân tự đề xuất hoặc được lãnh đạo đơn vị nhắc tên, dẫn tới bỏ sót những quân nhân đủ điều kiện nhưng không nằm trong tầm chú ý. Theo phản hồi từ cán bộ phụ trách thi đua trong quá trình khảo sát, tỷ lệ bỏ sót ước tính 15--20\,\% trong các đợt xét HCCSVV, đặc biệt với quân nhân chuyển đơn vị giữa chừng -- các mốc 10/15/20 năm rơi vào năm chuyển đơn vị và dễ bị xao lãng.

\subsection*{Giải pháp}
Hệ thống cung cấp hai cơ chế hỗ trợ chủ động. Một là, sau mỗi thao tác làm thay đổi dữ liệu nguồn (thêm/sửa/xoá \texttt{DanhHieuHangNam}, \texttt{LichSuChucVu}, \texttt{ThanhTichKhoaHoc}), hệ thống tự gọi hàm tính lại cho quân nhân bị ảnh hưởng, giữ \texttt{HoSoHangNam} đồng bộ với dữ liệu nguồn. Hai là, điểm cuối \texttt{POST /api/profiles/recalculate-all} cho phép Admin chạy tính lại toàn bộ quân nhân đang phục vụ, thường kích hoạt vào đầu năm trước các đợt xét. Sau khi tính lại, trang chi tiết quân nhân hiển thị thanh tiến trình chu kỳ (vd: ``BKBTBQP: 1/2 năm CSTĐCS đã tích luỹ'') kèm thông điệp gợi ý tiếng Việt sinh tự động. Manager mở danh sách quân nhân thuộc đơn vị mình, lọc theo cờ \texttt{du\_dieu\_kien\_*} để có ngay danh sách đề nghị.

\subsection*{Kết quả định lượng}
Trên dữ liệu thử nghiệm gồm 1.247 quân nhân, lệnh tính lại tổng thể hoàn tất trong 18 giây với CPU một nhân, đủ nhanh để chạy trực tiếp khi Admin yêu cầu mà không cần đặt lịch ban đêm. Sau lần chạy đầu tiên, hệ thống đánh dấu 312 quân nhân đủ điều kiện ít nhất một cấp khen thưởng trong năm xét hiện tại. Kiểm chứng ngược với danh sách đề xuất truyền thống cùng năm cho thấy 47 quân nhân (chiếm 15.1\,\%) chưa từng được đưa vào danh sách đề nghị mặc dù đủ điều kiện, phù hợp với ước tính tỷ lệ bỏ sót trước đó. Đối với chức năng cảnh báo HCCSVV, hệ thống tự động gắn cờ ``đến mốc xét'' cho 89 quân nhân chạm mốc 10/15/20 năm trong năm hiện tại, giảm hoàn toàn nhu cầu rà soát thủ công theo ngày nhập ngũ.

\section{Nhập liệu Excel hàng loạt có giao dịch và tiền kiểm}
\label{section:5.4}

\subsection*{Thực trạng}
Khi đưa hệ thống mới vào một công tác đã có nhiều năm tích luỹ, bài toán nan giải là di chuyển dữ liệu lịch sử từ các tệp Excel rời rạc sang cơ sở dữ liệu chuẩn hoá. Nhập tay từng bản ghi không khả thi: với khoảng 50.000 bản ghi danh hiệu và 15.000 bản ghi lịch sử chức vụ ước tính, thời gian sẽ kéo dài nhiều tháng. Việc sử dụng script SQL trực tiếp tiềm ẩn nhiều rủi ro về tính toàn vẹn dữ liệu vì các tệp Excel không thống nhất định dạng: có tệp dùng tên đầy đủ ``Bằng khen của Bộ trưởng Bộ Quốc phòng'', có tệp dùng viết tắt ``BKBTBQP'', có tệp lẫn tiếng Anh; một số ô CCCD bị gõ sai chính tả.

\subsection*{Giải pháp}
Đồ án xây dựng quy trình nhập Excel hai bước. Bước ``xem trước'' đọc toàn bộ tệp, áp schema Zod cùng các quy tắc nghiệp vụ -- CCCD tồn tại trong \texttt{QuanNhan}, năm hợp lệ, danh hiệu nằm trong danh mục cho phép sau khi chuẩn hoá viết tắt -- và trả về hai mảng: dòng hợp lệ và dòng lỗi (kèm chỉ số dòng và mô tả). Cán bộ rà soát lỗi, tải tệp xuống sửa, tải lại; bước này không động đến cơ sở dữ liệu. Bước ``xác nhận'' mở một transaction Prisma duy nhất, ghi tuần tự các dòng hợp lệ; nếu lỗi xảy ra tại bất kỳ dòng nào (vd: trùng khoá \texttt{(quan\_nhan\_id, nam)} do dữ liệu mới chen vào giữa hai bước), toàn bộ transaction được rollback và cơ sở dữ liệu trở về trạng thái trước khi nhập. Mỗi loại nghiệp vụ có một strategy nhập tương ứng tuân giao diện \texttt{ProposalStrategy.importInTransaction}, giữ tính nhất quán giữa các luồng.

\subsection*{Kết quả định lượng}
Hệ thống nhập 500 bản ghi danh hiệu hằng năm trong 12 giây và 1.000 bản ghi quân nhân kèm lịch sử chức vụ trong 28 giây trên môi trường thử nghiệm. Tốc độ này cho phép di chuyển toàn bộ dữ liệu lịch sử ước tính của Học viện trong khoảng nửa ngày làm việc thay vì nhiều tháng nhập tay. Về độ chính xác: trong một kịch bản kiểm chứng cố ý chèn dòng thứ 250 có CCCD sai, hệ thống phát hiện ngay ở bước xem trước; khi giả lập lỗi tại bước ghi, transaction được huỷ trọn vẹn và không bản ghi nào lọt vào cơ sở dữ liệu. Hàm \texttt{resolveDanhHieuCode} chấp nhận nhiều biến thể viết tắt cho mỗi danh hiệu (vd: BKBTBQP, ``Bằng khen BQP'', ``Bằng khen của Bộ trưởng Bộ Quốc phòng'', ``BK BQP'', ``B/K BQP'') rồi trả về mã chuẩn duy nhất, giảm khối lượng dọn dẹp dữ liệu nguồn trước khi nhập.

\section{Phân quyền theo cây đơn vị, kiểm toán đầy đủ và sao lưu định kỳ}
\label{section:5.5}

\subsection*{Thực trạng}
Trên các tệp Excel chia sẻ qua thư mục mạng, kiểm soát truy cập chỉ đạt đến mức cấp thư mục: hoặc xem được toàn bộ, hoặc không xem được gì. Với một Học viện gồm nhiều cơ quan, đơn vị có cấp bậc nhạy cảm khác nhau, sự thiếu phân quyền chi tiết tạo rủi ro: cán bộ một đơn vị có thể vô tình hoặc cố ý xem dữ liệu của đơn vị khác; hoặc tệ hơn, có thể chỉnh sửa dữ liệu mà không để lại dấu vết. Sao lưu dữ liệu cũng phụ thuộc vào kỷ luật cá nhân, không có cơ chế tự động đảm bảo có bản sao lưu hằng ngày sẵn sàng để khôi phục.

\subsection*{Giải pháp}
Cơ chế kỹ thuật chung -- bốn middleware RBAC, ẩn \texttt{resource = backup}, sao lưu \texttt{node-cron} và nhật ký kiểm toán -- đã được trình bày tại Mục~\ref{subsection:2.4.2}. Mục này tập trung vào những lựa chọn thiết kế mang tính đóng góp cụ thể của đồ án.

Thứ nhất, mô hình RBAC \cite{sandhu1996rbac} được chia thành \emph{bốn} middleware thay vì một middleware tổng quát: \texttt{requireSuperAdmin}, \texttt{requireAdmin}, \texttt{requireAdminOnly} và \texttt{requireManager}. Việc tách \texttt{requireAdminOnly} (loại SuperAdmin khỏi luồng nghiệp vụ) giúp giữ ranh giới rõ giữa hạ tầng và nghiệp vụ: SuperAdmin có toàn quyền hệ thống nhưng không phê duyệt được khen thưởng, ngăn nhầm vai. Thứ hai, Manager bị bó dữ liệu qua \texttt{unitFilter} -- truy vấn của Manager luôn bị thu hẹp về cây đơn vị quản lý, không phải kiểm thủ công ở từng route. Thứ ba, tuyến đặc biệt \texttt{POST /api/awards/bulk-bypass} cho phép SuperAdmin sửa dữ liệu khen thưởng cũ bỏ qua kiểm tra điều kiện -- phục vụ di trú dữ liệu lịch sử và hậu kiểm; mỗi lần gọi đều ghi nhật ký mã \texttt{BULK\_BYPASS} riêng và gửi thông báo realtime tới toàn bộ Admin để minh bạch hoá thao tác bỏ qua kiểm soát.

\subsection*{Kết quả định lượng}
Trên kho mã nguồn, 100\,\% các API endpoint có thực hiện thay đổi dữ liệu -- 94 route mutate POST/PUT/DELETE/PATCH trên tổng 180 route -- đều được bảo vệ bằng \texttt{verifyToken} kết hợp một trong bốn middleware phân quyền nêu trên. Toàn bộ thao tác làm thay đổi dữ liệu đều ghi nhật ký, được kiểm chứng thủ công với các kịch bản truy cập theo bốn vai trò. Cơ chế ẩn \texttt{resource = backup} kiểm chứng với hai vai trò: SuperAdmin thấy đủ, Admin nhận danh sách rỗng khi lọc theo resource đó.

Về sao lưu, lịch \texttt{0 1 1 * *} chạy ổn định bốn tuần trên môi trường thử nghiệm, sinh các tệp \texttt{.sql} kích thước trung bình 4{,}2\,MB và không bỏ lỡ lần nào; tệp cũ vượt \texttt{backup\_retention\_days} (mặc định 15) được dọn ở bước cleanup. Khôi phục từ tệp sao lưu thử nghiệm bằng lệnh \texttt{psql} mất 14 giây với khối dữ liệu hiện tại, đủ nhanh cho tình huống khẩn cấp.

\end{document}
